% TOP_PROBLEM: {"id": 3199, "title": "A generalization of Vizing's Theorem?", "category_id": 3, "category": {"id": 3, "name": "graph_theory", "display_name": "Graph Theory", "description": "Problems involving graphs, networks, and their properties.", "slug": "graph-theory", "order_index": 3, "created_at": "2026-07-31T15:26:25.671Z"}, "status": "open", "problem_number": "OPG-182", "set_id": 12}
% FIRST_POSTED: 2026-10-01
% ENTRY_KIND: statement_only
% UnsolvedMath ID 3199; source catalogue code OPG-182
% !TeX program = lualatex
% Definition and sources only; this file is not an LLM solution attempt.
\documentclass[11pt,a4paper]{article}
\usepackage[margin=25mm]{geometry}
\usepackage{fontspec}
\setmainfont{lmroman10-regular.otf}[BoldFont=lmroman10-bold.otf,
 ItalicFont=lmroman10-italic.otf,BoldItalicFont=lmroman10-bolditalic.otf]
\usepackage{amsmath,amssymb,mathtools}
\usepackage{unicode-math}
\setmathfont{latinmodern-math.otf}
\AtBeginDocument{\let\setminus\smallsetminus}
\usepackage{xurl,hyperref}
\hypersetup{colorlinks=true,urlcolor=blue}
\setcounter{secnumdepth}{0}
\setlength{\parindent}{0pt}
\setlength{\parskip}{5pt}
\setlength{\emergencystretch}{3em}
\begin{document}
\section*{A generalization of Vizing's Theorem?}\label{3199}
{\small UnsolvedMath ID 3199 · OPG-182 · Graph Theory · OpenGarden}
\subsection{Definitions and mathematical statement}
Fix integers $d\ge2$ and $r\ge1$. A finite simple $d$-uniform hypergraph $H=(V,\mathcal E)$ consists of a finite vertex set and a set of distinct $d$-element subsets of $V$, called edges. For a $(d-1)$-element set $S\subseteq V$, its codegree is the number of hyperedges containing $S$. Assume every such codegree is at most $r$.

The conjecture asks for an assignment
\[
 c:\mathcal E\longrightarrow\{1,\ldots,r+d-1\}
\]
such that $c(e)\ne c(f)$ whenever distinct hyperedges $e,f$ have $|e\cap f|=d-1$. This is a coloring of hyperedges with a specific conflict rule: pairs sharing fewer than $d-1$ vertices may receive the same color, even when their intersection is nonempty.

Equivalently, form the ordinary conflict graph with vertex set $\mathcal E$ and join exactly those pairs of hyperedges that overlap in $d-1$ vertices. The assertion is that this graph has chromatic number at most $r+d-1$ under the stated codegree bound. Both parameters are universal; the bound is independent of $|V|$ and $|\mathcal E|$.

When $d=2$, the hypergraph is an ordinary simple graph, codegree is vertex degree, and the conflict rule is the usual proper edge-coloring rule. This recovers the degree-plus-one edge-coloring bound. For larger $d$, replacing the conflict rule by the usual requirement that all intersecting hyperedges differ would change the problem substantially.
\subsection{Short English statement}
Can d-uniform hyperedges with codegree at most r be colored with r+d−1 colors when only overlaps of size d−1 create conflicts?
\subsection{Sources}
\begin{itemize}
\item[S1] ULAM.AI and the UnsolvedMath Contributors, supplied problems.json, record 3199 (OPG-182), OpenGarden. Catalogue identity and source transcription; CC BY 4.0.
\url{https://huggingface.co/datasets/ulamai/UnsolvedMath}.
\item[S2] Open Problem Garden, A generalization of Vizing's Theorem?. Original problem and discussion.
\url{http://www.openproblemgarden.org/op/a_generalization_of_vizings_theorem}.
\end{itemize}
\end{document}
